\documentclass[10pt,a4paper]{amsart}
\usepackage[margin=19mm]{geometry}
\usepackage{amsmath,amssymb,mathtools}
\usepackage[scr=rsfs,cal=euler]{mathalfa}
\usepackage{xfrac}
\usepackage[unicode,hidelinks,bookmarks=false]{hyperref}
\newcommand{\ie}{i.e.\ }
\newcommand{\cf}{cf.\ }
\newcommand{\resp}{respectively\ }
\newcommand{\Subsection}{\subsection}
\providecommand{\nobreakdash}{\nobreak\hbox{-}}
\setlength{\emergencystretch}{2em}
\multlinegap0pt
\usepackage{bm}
\usepackage{bbm}
\usepackage{dsfont}
\usepackage{xspace}
\usepackage{cancel}
\usepackage{mathtools}
\usepackage{amsthm}
\makeatletter
\@ifundefined{theorem}{\newtheorem{theorem}{Theorem}}{}
\@ifundefined{lemma}{\newtheorem{lemma}[theorem]{Lemma}}{}
\makeatother
\title[Property $(\diamond)$ for Ore extensions of small Krull dime]{Property $(\diamond)$ for Ore extensions of small Krull dimension}
\author{Ken Brown, Paula A.A.B. Carvalho, Jerzy Matczuk}
\date{Source: arXiv:2403.09239v1 (CC BY 4.0)}
\begin{document}
\maketitle

\noindent\emph{Source and licence.} Property $(\diamond)$ for Ore extensions of small Krull dimension. Version arXiv:2403.09239v1 (CC BY 4.0), distributed under the Creative Commons Attribution 4.0 International licence, as recorded in the arXiv licence metadata for this exact version and shown on the abstract page https://arxiv.org/abs/2403.09239v1. Full rights notice: licenses/arxiv-2403.09239-CC-BY-4.0.txt.\par\smallskip

\noindent\textit{Editorial scope.} Counted unit: the Lemma whose source label is \texttt{q algebraic} in the article (source0.tex lines 309--314, source label \texttt{q algebraic}), delivered together with the article's own complete original proof of it (source0.tex lines 316--577, one \texttt{proof} environment of 262 lines). No mathematics is written, repaired or re-derived: statement and proof are the article's own text line for line. Required context retained uncounted: irrepresentation (lines 305--306). Every definition, display and label of those lines is retained unchanged. Omitted from this package and not redistributed: 1--304, 307--308, 315--315, 578--1017. Nothing inside a retained excerpt is omitted or altered: every retained body is delivered exactly as the article's lines are written, so no omission receipt of any kind accompanies this package. Citations: the retained excerpts carry no citation command at all, so no bibliography entry of the article is transcribed and no reference is redistributed. Reference adaptation: none was needed and none was made; every reference inside the delivered unit resolves inside it, so no unresolved reference or citation is produced. Printed slips kept literal: no printed word, symbol, hypothesis or number of the counted statement or of its proof was altered. Layout and class: the article's own class is \texttt{article} with packages including amssymb, latexsym, ulem, amsrefs, amsmath, amsthm, color, marginnote, hyperref, cancel, a4wide; the wrapper is the lane's pinned \texttt{amsart} editorial wrapper at 10pt on A4, which restates the theorem-like environments the retained text needs and adds the article's own macro definitions that the retained text needs (none), with extra wrapper packages (bm, bbm, dsfont, xspace, cancel, mathtools). The delivered numbering is the unit's own. Assets: no retained span contains a figure environment, a graphics-inclusion command, a PDF-page inclusion, a table or a TikZ picture; no image file is copied into this package, the package declares no asset dependency and no image file is redistributed with it. Licence: version arXiv:2403.09239v1 of this article is distributed under the Creative Commons Attribution 4.0 International licence, as recorded in the arXiv licence metadata for this exact version and shown on the abstract page https://arxiv.org/abs/2403.09239v1. A full rights notice is delivered at licenses/arxiv-2403.09239-CC-BY-4.0.txt. Characters: every retained excerpt is pure ASCII, so the delivered engine, XeLaTeX with its default Latin Modern text font, reports no missing character.
 \begin{lemma}\label{irrepresentation} Keeping the above notation. Suppose that $f=f_Af_B$ is a $(g,\xi)$-irreducible  presentation of $f$. Then, for any $c\in C$ and $d\in D$, $d\ne c$ (as elements of the base field $k$).
 \end{lemma}

\begin{lemma}\label{q algebraic} Given $S$ as before suppose that
 \begin{equation}\label{E6}
(1+X+\theta+X\theta^2)(1-X+X\theta)=(1+Xt\theta)(a+b\theta+c\theta^2)
\end{equation}
for some  $t, a,b,c \in k[X]_{<X>}$. Then either $k$ has characteristic $2$ or $q$ is a root of unity.
\end{lemma}

\begin{proof}
Comparing the coefficients of the monomials in the equation (\ref{E6}) above we have that
\begin{equation}\label{E71}
1-X^2=a
\end{equation}
\begin{equation}\label{E72}
1-qX+X+X^2=Xt\alpha(a)+b
\end{equation}
\begin{equation}\label{E73}
qX+(1-q^2X)X=Xt\alpha(b)+c
\end{equation}
\begin{equation}\label{E74}
q^2X^2=Xt\alpha(c)
\end{equation}

Substituting $a$ from (\ref{E71}) in (\ref{E72}) we get
\begin{equation}	\label{E8}
1-qX+X+X^2=Xt(1-q^2X^2)+b
\end{equation}
and substituting $b$ from (\ref{E8}) in  (\ref{E73}) we obtain
\begin{equation}\label{E9}
c=qX+(1-q^2X)X-Xt[1-q^2X+qX+q^2X^2-qX\alpha(t)(1-q^4X^2)].
\end{equation}
Note that by (\ref{E73}) $c\in XR$, and by (\ref{E74}) $t$ is a unit in $R$. Write $c=Xc_1$, so that by (\ref{E74})
\begin{equation}\label{E10}
%q^2=tq\alpha(c_1)\\
c_1=q\alpha^{-1}(t^{-1}).
\end{equation}
From (\ref{E9}), (\ref{E10}) and applying $\alpha$ we get
\begin{equation}\label{reduced}
\begin{split}
q&=(1+q-q^3X)t\\&-(1+(1-q)q^2X+q^4X^2)t\alpha(t)\\&+q^2X(1-q^6X^2)t\alpha(t)\alpha^2(t).
\end{split}
\end{equation}

\noindent The problem is now reduced to deciding if there exists $t\in R$ that satisfies (\ref{reduced}). We can embed $R=k[X]_{<X>}$ in $k[[X]]$ and write $t=\displaystyle \sum_{i\geq 0} t_iX^i$ for $t_i \in k$. From (\ref{reduced}) it follows that
$$q=(1+q)t_0-t_0^2.$$
Hence
\begin{equation}\label{t0}t_0=1 \; \mbox{or}\; t_0=q.
\end{equation}
%and
%$$0=(1+q)t_1-q^3t_0-(t_1t_0+qt_0t_1+(1-q)q^2t_0^2)+q^2t_0^3.$$
%So either $t_0=q$.????????????????
\noindent Since $t\in k[X]_{<X>}$, we can write $t=\lambda\frac{f}{g}$ for some $\lambda\in k$, and  $f,g\in k[X]$ monic and co-prime such that $g\notin X k[X]$.  Let $f_0, g_0$ be respectively the constant terms of $f$ and $g$, so that by (\ref{t0})
\begin{equation}\label{degree0}
\lambda f_0= qg_0\quad \mbox{or}\quad \lambda f_0=g_0.
\end{equation}


\noindent Now (\ref{reduced}) can be written as
\begin{equation}\label{*}
\begin{split}
qg\alpha(g)\alpha^2(g)&=\left((1+q)-q^3X\right)\lambda f\alpha(g)\alpha^2(g)\\
&-\left(1+(1-q)q^2X+q^4X^2\right)\lambda^2f\alpha(f)\alpha^2(g)\\
&+q^2X(1-q^6X^2)\lambda^3f\alpha(f)\alpha^2(f).
\end{split}
\end{equation}%\footnote{We could use this eaquation to get (\ref{degree0}) but I have a problem with the $\sqrt{f_0}$}

%Consider the element $1+X+\theta+X\theta^2$ irreducibe in $S$ of type {\bf{(C)}}, \cite[Example 6.9]{BCM2} and $1-X+X\theta$ an irreducible element of type {\bf{(B)}}. We show that $v=(1+X+\theta+X\theta^2)(1-X+X\theta)$  can't be written as a product of an element of type {\bf{(B)}} by one of type {\bf{(C)}}. The proof is divided into 2 lemmas
\noindent Fix $m := \deg(g)$ and $n := \deg(f)$. If $m\leq n$ then in equation (\ref{*}) the LHS has degree $3m$ while the RHS has degree $3+3n$, a contradiction. If $m>n+1$, the LHS has degree $3m$ while the RHS has degree strictly less than $3m$. So $m=n+1$.

\medskip

Observe that if the conclusion of the lemma does not hold then it will not hold when $k$ is replaced by its algebraic closure. Hence, without of lost of  generality we may assume that $k$ is algebraically closed. Let  $Z_g:=\{z_1,\ldots, z_{n+1}\}$ denote the enumerated roots of $g$.  Since $f$ and $g$ are co-prime, by (\ref{*}) $f$ is a polynomial of degree $n$ such that every root of $f$ is a root either of $\alpha(g) $ or $\alpha^2(g) $. Thus the polynomials $f$ and $g$ satisfy the assumptions of Lemma \ref{irrepresentation} with $\xi=q^{-1}$. Fix the $(g,q^{-1})$-irreducible presentation $f=f_Af_B$ of $f$, so that $A$ and $B$ are  subsets of  $Z_g$ and  $f_A, f_B\in k[X]$ are monic polynomials  with the set $Z_{f_A}$ of all roots of  $f_A$ equal to $q^{-1}A$ and $Z_{f_B}=q^{-2}B$. Moreover the cardinality of  $C:=A\cap B$ is minimal for sets $A,B$ satisfying these conditions. Now we claim that 
\begin{equation}\label{Cempty} C\; =\; \emptyset.
\end{equation}

Write  $f_A=\prod_{z\in A}(X-q^{-1}z)$ and $f_B=\prod_{z\in B}(X-q^{-2}z)$. Note that

\begin{eqnarray}
\alpha(g) &=& q^{n+1}\prod_{i=1}^{n+1}(X-q^{-1}z_i)=q^{n+1}f_A\prod_{z\in B_0}(X-q^{-1}z)\prod_{z\in D}(X-q^{-1}z)\label{2fA},\\
\alpha^2(g) &=& q^{2n+2}f_B\prod_{z\in A_0}(X-q^{-2}z)\prod_{z\in D}(X-q^{-2}z).\label{2fB}
\end{eqnarray}

\noindent Combining (\ref{*}) with (\ref{2fA}) and (\ref{2fB}) we have (after cancelling $f_Af_B\prod_{z\in A_0}(X - q^{-2}z)$)

%\begin{equation}\label{2**}
%\begin{split}
%&q^{3n+4}g \prod_{z\in Z\backslash A}(x-q^{-1}z)\prod_{z\in Z\backslash B}(x-q^{-2}z)=\\
%&=\left((1+q)-q^3x\right)aq^{3n+3}f_A\prod_{z\in Z\backslash A}(x-q^{-1}z)f_B\prod_{z\in Z\backslash B}(x-q^{-2}z)\\
%&-\left(1+(1-q)q^2x+q^4x^2\right)a^2q^{3n+2}\prod_{z\in A}(x-q^{-2}z)\prod_{z\in B}(x-q^{-3}z)f_B\prod_{z\in Z\backslash B}(x-q^{-2}z)\\
%&-q^{n+8}x(x-\frac{1}{q^3})(x+\frac{1}{q^3})a^2\prod_{z\in A}(x-q^{-2}z)\prod_{z\in B}(x-q^{-3}z)\alpha^2(f)
%\end{split}
%\end{equation}

%Since
%\begin{equation}
%\prod_{z\in Z\backslash B}(x-q^{-2}z)=\prod_{z\in A_0}(x-q^{-2}z)\prod_{z\in D}(x-q^{-2}z)
%\end{equation}
%equation (\ref{2**}) can be written as (once divided by $\prod_{z\in A_0}(x-q^{-2}z))$)

\begin{equation}\label{2***}
\begin{split}
&q^{3n+4}g \prod_{z\in B_0}(X-q^{-1}z)\prod_{z\in D}(X-q^{-1}z)\prod_{z\in D}(X-q^{-2}z)=\\
&=\left((1+q)-q^3X\right)\lambda q^{3n+3}f_A\prod_{z\in B_0}(X-q^{-1}z)\prod_{z\in D}(X-q^{-1}z)f_B\prod_{z\in D}(X-q^{-2}z)\\
&-\left(1+(1-q)q^2X+q^4X^2\right)\lambda^2q^{3n+2}\prod_{z\in A}(X-q^{-2}z)\prod_{z\in B}(X-q^{-3}z)f_B\prod_{z\in D}(X-q^{-2}z)\\
&-q^{n+8}X(X-\frac{1}{q^3})(X+\frac{1}{q^3})\lambda^3\prod_{z\in C}(X-q^{-2}z)\prod_{z\in B}(X-q^{-3}z)\alpha^2(f)
\end{split}
\end{equation}

Now $\prod_{z\in C}(X-q^{-2}z)$ divides $f_B$ since $C\subseteq B$, so $\prod_{z\in C}(X-q^{-2}z)$ divides the RHS of equation (\ref{2***}). Since for each $z\in C, (X-q^{-2}z)$ divides $f$, it does not divide $g$. Also, by Lemma \ref{irrepresentation},  for each $z\in C, (X-q^{-2}z)$ does not divide $\prod_{z\in D}(X-q^{-2}z)$. It therefore follows that
\begin{equation}\label{division}\prod_{z\in C}(X-q^{-2}z) \textit{ divides } \prod_{z\in B_0}(X-q^{-1}z)\prod_{z\in D}(X-q^{-1}z).\end{equation}

\noindent Hence, given any $z\in C$, either there is a $b\in B_0$ such that $q^{-2}z=q^{-1}b$ or there is a $d\in D$ such that $q^{-2}z=q^{-1}d$. In both these cases $q^{-1}z=b$ or $q^{-1}z=d$ it follows that  $q^{-1}z$  is a root of $f_A$ and of $g$, contradicting the fact that $f$ and $g$ are coprime. So $C=\emptyset$ - that is,  (\ref{Cempty}) is proved.

\medskip

Now we are able to complete the proof of the lemma. Without lost of generality assume that $D=Z_g\setminus (A\cup B)=\{z_{n+1}\}$. By (\ref{Cempty}), we can now rewrite (\ref{2***}) as
%\bigskip


%Let $Z_g=\{z_1,\ldots,z_{n+1}\}$ be the set of zeros of  $g=\prod_{i=1}^n(x-z_i)$, $f=af_Af_B$ for $f_A=\prod_{z\in A}(x-q^{-1}z)$ and $f_B=\prod_{z\in B}(x-q^{-2}z)$ for $A,B\subseteq Z_g$ such that $A\cap B=\emptyset$.

%Without loss of generality assume that $z_{n+1} \notin A\cup B$. Note that

%\begin{eqnarray}
%\alpha(g) &=& q^{n+1}\prod_{i=1}^{n+1}(x-q^{-1}z_i)=q^{n+1}f_A\prod_{z\in Z\backslash A}(x-q^{-1}z)\label{fA}\\
%\alpha^2(g) &=& q^{2n+2}f_B\prod_{z\in Z\backslash B}(x-q^{-2}z)\label{fB}
%\end{eqnarray}

%Combining (\ref{*}) with (\ref{fA}) and (\ref{fB}) we have

%\begin{equation}\label{**}
%\begin{split}
%&q^{3n+4}g \prod_{z\in Z\backslash A}(x-q^{-1}z)\prod_{z\in Z\backslash B}(x-q^{-2}z)=\\
%&=\left((1+q)-q^3x\right)aq^{3n+3}f_A\prod_{z\in Z\backslash A}(x-q^{-1}z)f_B\prod_{z\in Z\backslash B}(x-q^{-2}z)\\
%&-\left(1+(1-q)q^2x+q^4x^2\right)a^2q^{3n+2}\prod_{z\in A}(x-q^{-2}z)\prod_{z\in B}(x-q^{-3}z)f_B\prod_{z\in Z\backslash B}(x-q^{-2}z)\\
%&+q^{n+6}x(x-\frac{1}{q^3})(x+\frac{1}{q^3})a^2\prod_{z\in A}(x-q^{-2}z)\prod_{z\in B}(x-q^{-3}z)\alpha^2(f)
%\end{split}
%\end{equation}

%Since
%\begin{equation}
%\prod_{z\in Z\backslash B}(x-q^{-2}z)=\prod_{z\in A}(x-q^{-2}z)(x-q^{-2}z_{n+1})
%\end{equation}


\begin{equation}\label{***}
\begin{split}
&q^{3n+4}g \prod_{z\in Z_g\backslash A}(X-q^{-1}z)(X-q^{-2}z_{n+1})=\\
&=\left((1+q)-q^3X\right)\lambda q^{3n+3}f_A\prod_{z\in Z_g\backslash A}(X-q^{-1}z)f_B(X-q^{-2}z_{n+1})\\
&-\left(1+(1-q)q^2X+q^4X^2\right)\lambda^2q^{3n+2}\prod_{z\in A}(X-q^{-2}z)\prod_{z\in B}(X-q^{-3}z)f_B(X-q^{-2}z_{n+1})\\
&-q^{n+8}X(X-\frac{1}{q^3})(X+\frac{1}{q^3})\lambda^3\prod_{z\in B}(X-q^{-3}z)\alpha^2(f).
\end{split}
\end{equation}


%below killed on Nov the 2nd to preplace the part Ken had some issues with
%Since $f$ and $g$ are coprimes, so are $\alpha^2(f)$ and $\alpha^2(g)$ hence $(X-q^{-2}z_{n+1})$ does not divides $\alpha^2(f)$. From equation (\ref{***}) we have
%\begin{equation}\label{zn}z_{n+1}=\pm \frac{1}{q}\quad \vee\quad \exists b_1\in B: z_{n+1}=q^{-1}b_1. \end{equation}
%If the first case doesn't happen, equation (\ref{***}) can be written as
%\begin{equation}\label{simp_1}
%\begin{split}
%&q^{3n+4}g \prod_{z\in Z_g\backslash A}(X-q^{-1}z)=\\
%&=\left((1+q)-q^3X\right)\lambda q^{3n+3}f_A\prod_{z\in Z_g\backslash A}(X-q^{-1}z)f_B\\
%&-\left(1+(1-q)q^2X+q^4X^2\right)\lambda^2q^{3n+2}\prod_{z\in A}(X-q^{-2}z)\prod_{z\in B}(X-q^{-3}z)f_B\\
%&-q^{n+8}X(X-\frac{1}{q^3})(X+\frac{1}{q^3})\lambda^3\prod_{z\in B\backslash\{b_1\}}(X-q^{-3}z)\alpha^2(f)
%\end{split}
%\end{equation}
%Note that $Z_g\backslash A=B\cup \{z_{n+1}\}$ and that $(X-q^{-1}z_{n+1})$ is a factor of the LHS of equation (\ref{simp_1}), $(X-q^{-1}z_{n+1})=(X-q^{-2}b_1)$ divides $f_B$, $\alpha^2(g)$ and doesn't divide $\alpha^2(f)$, so
%$$z_{n+1}=\pm \frac{1}{q^2}\quad \vee\quad \exists b_2\in B\backslash\{b_1\}: z_{n+1}=q^{-1}b_1=q^{-2}b_2.$$
%Once again, if the first case doesn't hold, dividing by $(X-q^{-1}z_{n+1})=(X-q^{-2}b_1)=(X-q^{-3}b_2)$ we have
%\begin{equation}\label{simp_2}
%\begin{split}
%&q^{3n+4}g \prod_{z\in B}(X-q^{-1}z)=\\
%&=\left((1+q)-q^3X\right)\lambda q^{3n+3}f_A\prod_{z\in Z\backslash A}(X-q^{-1}z)\prod_{z\in B\backslash\{b_1\}}(X-q^{-2}z)\\
%&-\left(1+(1-q)q^2X+q^4X^2\right)\lambda^2q^{3n+2}\prod_{z\in A}(X-q^{-2}z)\prod_{z\in B}(X-q^{-3}z)\prod_{z\in B\backslash\{b_1\}}(X-q^{-2}z)\\
%&-q^{n+8}X(X-\frac{1}{q^3})(X+\frac{1}{q^3})\lambda^3\prod_{z\in B\backslash\{b_1,b_2\}}(X-q^{-3}z)\alpha^2(f)
%\end{split}
%\end{equation}
%and $(X-q^{-1}b_1)=(X-q^{-2}b_2)$ is a factor of the LHS of equation (\ref{simp_2}), of $\prod_{z\in B\backslash\{b_1\}}(X-q^{-2}z)$  and not of $\alpha^2(f)$, hence
%$$z_{n+1}=\pm \frac{1}{q^3}\quad \vee\quad \exists b_3\in B\backslash\{b_1, b_2\}: z_{n+1}=q^{-1}b_1=q^{-2}b_2=q^{-3}b_3.$$
%comparing the number of factors of $\prod_{z\in B}(X-q^{-1}z)$  and of $\prod_{z\in B\backslash\{b_1,b_2\}}(X-q^{-3}z)$ allow us to conclude that
%either $z_{n+1}=\pm\frac{1}{q}$ or there is $1\leq l\leq |B|$ such that $z_{n+1}=\pm\frac{1}{q^{l+1}}.$ So there is $0\leq l\leq |B|$ such that $z_{n+1}=\pm\frac{1}{q^{l+1}}.$
% there will be a $l\leq |B|+1$ such that $b_l=\pm \frac{1}{q}$ and $z_{n+1}=q^{-l}b_l=\pm\frac{1}{q^{l+1}}.$

\noindent We claim that 
\begin{equation}\label{ellclaim} \exists \, \ell \textit{ such that } 0\leq \ell \leq |B| \textit{ and } z_{n+1}=\pm q^{-\ell -1}.
\end{equation}

Since $f$ and $g$ are coprime, so are $\alpha^2(f)$ and $\alpha^2(g)$.  Hence $(X-q^{-2}z_{n+1})$ does not divide $\alpha^2(f)$. Also $z_{n+1}\neq 0$, as $g\notin Xk[X]$,  and from equation (\ref{***}) we have
\begin{equation}\label{zn}z_{n+1}=\pm \frac{1}{q}\quad \vee\quad \exists b_1\in B: z_{n+1}=q^{-1}b_1. \end{equation}
If the first case happens then (\ref{ellclaim}) follows with $l=0$. Suppose now that the second case of (\ref{zn}) holds. Then equation (\ref{***}) can be written as
\begin{equation}\label{simp_1}
\begin{split}
&q^{3n+4}g \prod_{z\in Z_g\backslash A}(X-q^{-1}z)=\\
&=\left((1+q)-q^3X\right)\lambda q^{3n+3}f_A\prod_{z\in Z_g\backslash A}(X-q^{-1}z)f_B\\
&-\left(1+(1-q)q^2X+q^4X^2\right)\lambda^2q^{3n+2}\prod_{z\in A}(X-q^{-2}z)\prod_{z\in B}(X-q^{-3}z)f_B\\
&-q^{n+8}X(X-\frac{1}{q^3})(X+\frac{1}{q^3})\lambda^3\prod_{z\in B\backslash\{b_1\}}(X-q^{-3}z)\alpha^2(f)
\end{split}
\end{equation}

Let $b'_1,\ldots,b'_{|B|}$ be an enumeration of the roots of $g$ that are in $B$, assuming without loss of generality that 
\begin{equation}\label{add1} b'_1 \; = \; qz_{n+1} \;= \; b_1.
\end{equation}
(Note that we can have multiple roots, hence it may happen that $b'_i=b'_j$ as elements of $k$ for $i\neq j$). If $q^2z_{n+1}$ is one of the elements of the list $b'_2,\ldots, b'_{|B|}$ we let $b_2=q^2z_{n+1}$. Without loss of generality we can assume that $b_2=b'_2$. We proceed with $q^3z_{n+1}$, if it is in the list $b'_3,\ldots,b'_{|B|}$ then we write $b_3=q^3z_{n+1}$, reordering the roots in the list $b'_3,\ldots,b'_{|B|}$ we may assume that $b_3=b'_3$. Iterate the procedure until we reach $l$ such that $b_i=q^{i}z_{n+1}=b'_i$ for all $i\leq l$ but $q^{l+1}z_{n+1}$ is not in the list $b'_{l+1},\ldots, b'_{|B|}$. Note that $l\leq |B|$ and that
$$q^{-1}z_{n+1}=q^{-2}b_1=q^{-3}b_2,$$
$$ q^{-1}b_i=q^{-2}b_{i+1}=q^{-3}b_{i+2}, \forall i\in\{1,\ldots, l-2\},$$
$$q^{-1}b_{l-1}=q^{-2}b_l.$$
Now we have by (\ref{add1}) that
$$\prod_{z\in\{b_1,\ldots,b_{l-2}\}\cup\{z_{n+1}\}}(X-q^{-1}z)=\prod_{z\in\{b_1,\ldots,b_{l-1}\}}(X-q^{-2}z)=\prod_{z\in\{b_2,\ldots,b_{l}\}}(X-q^{-3}z).$$

So, from (\ref{simp_1}) and after cancelling the factor above, we get
\begin{equation}\label{simp_1*}
\begin{split}
&q^{3n+4}g \prod_{z\in B\backslash \{b_1,\ldots,b_{l-2}\}}(X-q^{-1}z)=\\
&=\left((1+q)-q^3X\right)\lambda q^{3n+3}f_A\prod_{z\in Z_g\backslash A}(X-q^{-1}z)\prod_{z\in B\backslash \{b_1,\ldots,b_{l-1}\}}(X-q^{-2}z)\\
&-\left(1+(1-q)q^2X+q^4X^2\right)\lambda^2q^{3n+2}\prod_{z\in A}(X-q^{-2}z)\prod_{z\in B}(X-q^{-3}z)\prod_{z\in B\backslash \{b_1,\ldots,b_{l-1}\}}(X-q^{-2}z)\\
&-q^{n+8}X(X-\frac{1}{q^3})(X+\frac{1}{q^3})\lambda^3\prod_{z\in B\backslash\{b_1,\ldots,b_l\}}(X-q^{-3}z)\alpha^2(f)
\end{split}
\end{equation}

Now note that $(X-q^{-1}b_{l-1})$ divides the LHS and is equal to $(X-q^{-2}b_l)$, which divides $\prod_{z\in B\backslash \{b_1,\ldots,b_{l-1}\}}(X-q^{-2}z)$. Since $q^{-2}b_l$ is a root of $\alpha^2(g)$ can't be a root of $\alpha^2(f)$, neither can it be $0$. Hence either $q^{-2}b_l=q^{-3}b$ for some $b\in B\backslash\{b_1,\ldots,b_l\}$ or $q^{-2}b_l=\pm q^{-3}$. Notice that the first equation cannot hold as otherwise 
$$q^{l+1}z_{n+1}=qb_l=b\in B\backslash\{b_1,\ldots,b_l\},$$ 
which is impossible because of the definition of $l$. Thus $q^{-2}b_l=\pm q^{-3}$  and consequently $z_{n+1}=q^{-l}b_l=\pm q^{-l-1}$, proving (\ref{ellclaim}).


%%%%%%%%%%%%%%%%%%%%%%%%%%%%%%%%%%%%%%%%%%%%%%%%%%
%%%%%%%%%%%%%%%%%%%%%%%%%%%%%%%%%%%%%%%%%%%%%%%%%%

By equation (\ref{degree0}), we have

\begin{equation}
\lambda q^{-|A|-2|B|}\prod_{i=1}^nz_i=-q\prod_{i=1}^{n+1}z_i\quad \mbox{or}\quad \lambda q^{-|A|-2|B|}\prod_{i=1}^nz_i=-\prod_{i=1}^{n+1}z_i
\end{equation}
hence
\begin{equation}\label{a}
\lambda =q^{|A|+2|B|+1}z_{n+1}=\pm q^{|A|+2|B|-l}=\pm q^{n+|B|-l}\quad \mbox{or}\quad \lambda=\pm q^{n+|B|-l-1}
\end{equation}
for some $0\leq l\leq |B|$.

Looking at the leading coefficient of equation (\ref{*}), we have
$$q^{3n+4}=-q^{3n+6}\lambda-q^{3n+6}\lambda^2-q^{3n+8}\lambda^3$$
so $$0=1+q^2\lambda+q^2\lambda^2+q^4\lambda^3=(\lambda+\frac{1}{q^2})(q^4\lambda^2+q^2)$$
from what follows that either $\lambda=-q^{-2}$ or $\lambda^2=-q^{-2}$
and by (\ref{a}) we have
\begin{equation}
1=\pm q^{n+|B|-l+2}\quad \mbox{or}\quad 1=\pm q^{n+|B|-l+1}
\end{equation}
or
\begin{equation}\label{final}
-1=q^{2n+2|B|-2l+2}\quad \mbox{or}\quad -1=q^{2n+2|B|-2l}
\end{equation}
By the choice of $l$, $n+|B|-l+1\geq n+1$ and if $\mathrm{char}(k)\neq 2$ $q$ is a root of unity.
%Note that if  $n+|B|-l=0$ then $n=0=|B|$ so $f=1$ and $g=(X-z)$.

%In characteristic 2, equation (\ref{*}), implies that  $z=1/q$ and we get
%\begin{equation}\label{char2}
%\begin{split}
%q(X-z)(qX-z)q^2&=\left((1+q)-q^3X\right)a(qX-z)q^2\\
%&-\left(1+(1-q)q^2X+q^4X^2\right)a^2q^2\\
%&+q^8X(X+\frac{1}{q^3})a^3
%\end{split}
%\end{equation}
%but then $\frac{1}{q^2}$ is a root of
%$$-\left(1+(1-q)q^2X+q^4X^2\right)a^2q^2+q^8X(X+\frac{1}{q^3})a^3$$
%so $(1+q)+q(1+q)a=0$ and either $q=1$ or $a=\frac{1}{q}$.


%This shows that $q$ is a root of unity.
% $q$ is algebraic over the prime subfield of $k$, so the thesis follows.
\end{proof}

\end{document}
